\begin{question}
\noindent In a school science lab, a microscope projects a slide image onto a screen. The lens system inverts the image and enlarges it by scale factor $-2$ about a fixed point $O$ on the screen (the optical centre).

\noindent Triangle $PQR$ represents a specimen shown on the slide, with $O$ marking the optical centre of the lens. The coordinates are shown on the grid below.

\begin{center}
\begin{tikzpicture}[scale=0.55, transform shape]
    \definecolor{borderColour}{HTML}{000000}
    \definecolor{fillColour}{HTML}{D8D8D8}

    \def\axisLimit{6}

    % Draw grid
    \draw[step=1cm,gray,very thin] (-\axisLimit,-\axisLimit) grid (\axisLimit,\axisLimit);

    % Draw axes
    \draw[thick,->] (-\axisLimit-0.5,0) -- (\axisLimit+0.5,0) node[anchor=north] {$x$};
    \draw[thick,->] (0,-\axisLimit-0.5) -- (0,\axisLimit+0.5) node[anchor=east] {$y$};

    % Label x-axis and y-axis
    \foreach \x in {-\axisLimit,...,-1,1,2,...,\axisLimit}
        \draw (\x cm,1pt) -- (\x cm,-1pt) node[anchor=north] {$\x$};
    \foreach \y in {-\axisLimit,...,-1,1,2,...,\axisLimit}
        \draw (1pt,\y cm) -- (-1pt,\y cm) node[anchor=east] {$\y$};

    % Centre of enlargement O
    \coordinate (O) at (-1,1);
    \draw[red, thick] (O) node[cross out, draw=red, inner sep=2pt, outer sep=0pt] {};
    \node at (-1.5,1.5) {$O$};

    % Original triangle PQR
    \coordinate (P) at (0,0);
    \coordinate (Q) at (0,-1);
    \coordinate (R) at (-2,0);

    \draw[fill=fillColour, draw=borderColour, line width=1pt] (P) -- (Q) -- (R) -- cycle;
    \node at (0.3,0.3) {$P$};
    \node at (0.3,-1.3) {$Q$};
    \node at (-2.4,0.3) {$R$};

\end{tikzpicture}
\end{center}

\noindent (a) On the grid, draw the image of triangle $PQR$ after an enlargement of scale factor $-2$, centre $O(-1,1)$. Label the image $P'Q'R'$.

\vspace{5cm}

\noindent (b) State the coordinates of $P'$, $Q'$ and $R'$.

\vspace{3cm}

\begin{flushright}
\answerplain{4}
\end{flushright}
\end{question}
